激光雷达点云技术因其在高精度三维重建中的广泛应用而备受关注，尤其是在自动驾驶、机器人导航和文化遗产保护等领域。随着点云数据采集技术的进步，如何高效、准确地从点云数据中重建三维模型成为研究热点。本文针对基于激光雷达点云的高精度三维重建技术展开研究，旨在探索传统算法与新兴隐式重建方法的性能差异及其在实际应用中的表现。

首先，本文分析了点云三维重建的基本理论与算法框架，深入研究了传统算法Alpha Shapes和Ball Pivoting Algorithm（BPA）滚球算法。通过构建三维显示表面重建算法，从多种开源三维点云数据集中重建出三维表面，并采用多种量化指标（倒角距离，豪斯多夫距离）对比不同传统算法的性能，验证了传统方法在三维点云数据中的有效性。

其次，本文重点研究了基于神经符号距离函数的隐式重建方法StEik，实现了一种三维隐式表面重建算法。实验结果表明，StEik算法在处理复杂点云数据时，能够显著提升重建精度和效率，尤其在噪声鲁棒性和细节还原方面优于传统方法，证明了其在高精度重建中的优越性。

为进一步验证算法的实际应用能力，本文采用3D打印技术制作多种物体，并利用三维扫描仪采集精确点云数据集。通过对比不同算法在自采集数据集上的性能，发现StEik算法在复杂场景中表现出更高的稳定性和适应性，为实际工程应用提供了可靠的技术支持。

综上，本文研究了传统与隐式方法的三维重建技术方案，通过理论分析和实验验证，为高精度三维重建领域提供了研究思路和技术路径，对推动激光雷达点云重建技术的实际应用具有重要意义。
